\documentclass[../../../include/open-logic-section]{subfiles}

\begin{document}

\olfileid[hi]{his}{set}{cantorplane}
\olsection{रेखा और समतल पर कांटोर}

श्रेडर--बर्नस्टाइन के प्रमाण से जुड़ी कुछ परिस्थितियाँ उस इतिहास से संबंधित
हैं जिसकी चर्चा हमने \olref[his][set][pathology]{sec} में की थी। याद करें
कि 1877 में कांटोर ने सिद्ध किया कि किसी वर्ग पर उतने ही बिंदु हैं जितने उसकी
किसी एक भुजा पर। यहाँ हम उनका (पहला प्रयास किया हुआ) प्रमाण प्रस्तुत करेंगे।

मान लें $\unitline$ इकाई रेखा है, अर्थात् बिंदुओं का समुच्चय $[0,1]$। मान
लें $\unitsquare$ इकाई वर्ग है, अर्थात् बिंदुओं का समुच्चय $\unitline
\times \unitline$। इन पदों में कांटोर ने सिद्ध किया कि
$\cardeq{\unitline}{\unitsquare}$। उन्होंने डेडेकिंड को एक टिप्पणी लिखी,
जिसमें मूलतः निम्नलिखित तर्क था।

\begin{thm}\ollabel{thm:cantorplane}
$\cardeq{\unitline}{\unitsquare}$
\end{thm}

\begin{proof}[प्रमाण: पहला भाग।]
$a, b \in \unitline$ नियत करें। उन्हें द्विआधारी संकेतन में लिखें, ताकि
हमारे पास $0$ और $1$ के अनंत अनुक्रम $a_1$, $a_2$, \dots, तथा $b_1$,
$b_2$, \dots हों, ऐसे कि:
\begin{align*}
a &= 0.a_1a_2a_3a_4\dots\\
b &= 0.b_1b_2b_3b_4\dots
\intertext{अब निम्नलिखित फलन $f \colon \unitsquare \to \unitline$ पर विचार करें} 
f(a, b) & = 0.a_1b_1a_2b_2a_3b_3a_4b_4\dots
\end{align*}
अब $f$ !!a{injection} है, क्योंकि यदि $f(a, b) = f(c,d)$, तो सभी $n \in
\Nat$ के लिए $a_n = c_n$ और $b_n = d_n$, अतः $a = c$ और $b = d$।
\end{proof}

दुर्भाग्य से, जैसा डेडेकिंड ने कांटोर को बताया, यह मूल प्रश्न का उत्तर नहीं
देता। $0.\dot{1}\dot{0} = 0.1010101010\ldots$ पर विचार करें। हमें चाहिए कि
$f(a,b) = 0.\dot{1}\dot{0}$, जहाँ:
\begin{align*}
a&= 0.\dot{1}\dot{1} = 0.111111\ldots\\
b&= 0
\end{align*}
किंतु $a = 0.\dot{1}\dot{1} = 1$। इसलिए जब हम कहते हैं ``$a$ और $b$ को
द्विआधारी संकेतन में लिखें'', तो हमें चुनना होगा कि \emph{कौन-सा} संकेतन
प्रयोग करें; और चूँकि $f$ को एक \emph{फलन} होना है, हम दो संभावित संकेतनों
में से केवल \emph{एक} प्रयोग कर सकते हैं। किंतु यदि, उदाहरणतः, हम सरल
संकेतन प्रयोग करके $a$ को ``$1.000\ldots$'' लिखें, तो हमारे पास ऐसा कोई युग्म
$\tuple{a, b}$ नहीं जिसके लिए $f(a, b) = 0.\dot{1}\dot{0}$।

संक्षेप में: डेडेकिंड ने बताया कि कुछ आवर्ती दशमलव प्रसारों की संभावना के
कारण कांटोर का फलन $f$ !!a{injection} तो है, किंतु !!a{surjection}
\emph{नहीं}। अतः कांटोर ने केवल $\cardle{\unitsquare}{\unitline}$ दिखाया
है, $\cardeq{\unitsquare}{\unitline}$ \emph{नहीं}।

कांटोर ने लगभग तत्काल डेडेकिंड को उत्तर लिखकर मूलतः सुझाया कि प्रमाण को इस
प्रकार पूरा किया जा सकता है:

\begin{proof}[प्रमाण: पूर्ण।]
अतः हमने दिखाया है कि $\cardle{\unitsquare}{\unitline}$। किंतु स्पष्टतः
$\unitline$ से $\unitsquare$ में !!a{injection} है: रेखा को वर्ग की किसी
एक भुजा पर समतल रख दें। अतः $\cardle{\unitline}{\unitsquare}$ और
$\cardle{\unitsquare}{\unitline}$। श्रेडर--बर्नस्टाइन
(\olref[sfr][siz][sb]{thm:schroder-bernstein}) से,
$\cardeq{\unitline}{\unitsquare}$।
\end{proof}

किंतु निस्संदेह कांटोर अंतिम पंक्ति को इन पदों में पूरा नहीं कर सकते थे,
क्योंकि श्रेडर--बर्नस्टाइन प्रमेय तब तक सिद्ध नहीं हुआ था। वास्तव में, यद्यपि
कांटोर ने बाद में इसे एक सामान्य अनुमान के रूप में सूत्रबद्ध किया, फिर भी
1897 तक इसका संतोषजनक प्रमाण नहीं हुआ। (और इसलिए 1877 में बाद में कांटोर ने
\olref{thm:cantorplane} का एक भिन्न प्रमाण दिया, जो श्रेडर--बर्नस्टाइन से
होकर नहीं जाता था।)

\end{document}
